\section{Physics Layer I: Landscapes, Barriers, and Energy-Style Intuition}

Physics language remains useful when used carefully. Energy-style intuition is
helpful because it compresses several qualitative ideas:
\begin{itemize}[nosep]
  \item deep wells are hard to leave;
  \item shallow wells are easy to destabilize;
  \item nearby ridges or saddles determine whether a small push can redirect
        the path;
  \item noise can produce barrier crossing even without perfect control.
\end{itemize}

\begin{discussion}
The pedagogical benefit of landscape language is not that the system has a
literal mechanical energy known to the reader. The benefit is that it provides
a tractable mental model for persistence and escape. In this book, the energy
picture should therefore be treated as a structured analogy supported by
medium-lane mathematical reasoning and by the strong empirical permission to
talk about stability and transition in biology.
\end{discussion}
